% TOP_PROBLEM: {"id": 30006996, "title": "Foulkes inclusion for the sixth symmetric power", "category_id": 4, "category": {"id": 4, "name": "algebra", "display_name": "Algebra", "description": "Group theory, ring theory, field theory, and algebraic structures.", "slug": "algebra", "order_index": 4, "created_at": "2026-07-31T15:26:25.671Z"}, "status": "open"}
% FIRST_POSTED: 2026-10-06
% ATTEMPT_STATUS: solved
% OPENAI_CLAIM: {"schema_version":1,"source_repo":"https://github.com/openai/math","source_commit":"adc7f1241b42e322a6451854ab7e4b4c146bf78a","release_date":"2026-10-06","match":"full","resolution":"proved","independently_reviewed":false,"second_pass":"agreed","families":[{"family":"210","title":"Foulkes' conjecture for sixth powers and quadratic stabilization","subject":"Algebra","manuscripts":[{"title":"Foulkes' conjecture for the sixth symmetric power","pdf_path":"preprints/Foulkes-Conjecture-for-the-Sixth-Symmetric-Power-September-25-2026/main.pdf","readme_path":"preprints/Foulkes-Conjecture-for-the-Sixth-Symmetric-Power-September-25-2026/README.md","date":"2026-09-25","theorem_ref":"thm:main (sections/introduction.tex:19)","statement_sources":[{"path":"preprints/Foulkes-Conjecture-for-the-Sixth-Symmetric-Power-September-25-2026/build/sections/introduction.tex","line":19,"label":"thm:main"}]}],"lean":{"doc_path":"lean/docs/210.md","comparators":[],"declarations":[],"covers_main_theorem":false},"reasoning_trace":null}]}
\documentclass[11pt]{article}
\usepackage[margin=25mm]{geometry}
\usepackage{fontspec}
\setmainfont{Latin Modern Roman}
\usepackage{amsmath,amssymb,mathtools}
\usepackage{unicode-math}
\setmathfont{Latin Modern Math}
\usepackage{enumitem}
\usepackage{xurl,hyperref,bookmark}
\hypersetup{colorlinks=true,urlcolor=blue}
\setcounter{secnumdepth}{0}
\setlength{\parindent}{0pt}
\setlength{\parskip}{5pt}
\setlength{\emergencystretch}{3em}
\begin{document}
\section*{Foulkes inclusion for the sixth symmetric power}

\subsection{OpenAI's claim}
OpenAI presents a proof in the selected manuscript. The sixth-power Foulkes comparison is claimed for every second exponent $b\geq6$ and every finite complex dimension.

\subsection{Mathematical statement of the claimed result}
\paragraph{Result 1.} For each integer \(b\ge6\) and finite-dimensional complex \(V\), the source claims a \(\operatorname{GL}(V)\)-equivariant inclusion
\[
 \operatorname{Sym}^6(\operatorname{Sym}^bV)
 \hookrightarrow\operatorname{Sym}^b(\operatorname{Sym}^6V).
\]
Equivalently, \(h_b[h_6]-h_6[h_b]\) has nonnegative Schur coefficients, where brackets denote plethysm and \(h_j\) the complete symmetric function.

\subsection{Outline of the argument}
Large degrees follow from canonical-map stabilization. For the remaining degrees the proof compares source character bounds with independently amplified target multiplicities, uses highest-weight shifts to reduce partitions, and finishes a narrow residual band by exact rational-function certificates.

\subsection{Scope relative to this problem}
The primary theorem gives the requested equivariant injection for every b>=6 and every finite-dimensional complex V, including 6<=b<30. lean/docs/210.md explicitly restricts the formal sixth-power range to b>=30 and excludes this full companion range. Within these assumptions, this is a claimed affirmative answer to the stated question.

\subsection{Formal verification}
For family 210, the formal scope is canonical surjectivity in the quadratic range, its $a=1$ case, and reverse embeddings. For sixth powers this only covers $b\geq30$, rather than every $b\geq6$. No Comparator is assigned to the selected manuscript here; the family page selects other papers or supporting statements. This site has not executed the Lean code. The scope descriptions do not certify the remaining manuscript arguments.

\subsection{Status caveat}
OpenAI attributes these manuscripts to work by an internal model. Its release cautions that arguments without formalization can contain mistakes. This catalogue has not independently checked the proofs; the labels report the mathematical claims attributed to OpenAI.

\subsection{Sources}
\begin{itemize}
\item OpenAI, \emph{Foulkes' conjecture for the sixth symmetric power}, 2026-09-25; citation key \texttt{OAI:Foulkes-Conjecture-for-the-Sixth-Symmetric-Power-September-25-2026}. \href{https://github.com/openai/math/blob/adc7f1241b42e322a6451854ab7e4b4c146bf78a/preprints/Foulkes-Conjecture-for-the-Sixth-Symmetric-Power-September-25-2026/main.pdf}{PDF (GitHub)}; \href{https://github.com/openai/math/raw/adc7f1241b42e322a6451854ab7e4b4c146bf78a/preprints/Foulkes-Conjecture-for-the-Sixth-Symmetric-Power-September-25-2026/main.pdf}{PDF (direct)}; \href{https://github.com/openai/math/blob/adc7f1241b42e322a6451854ab7e4b4c146bf78a/preprints/Foulkes-Conjecture-for-the-Sixth-Symmetric-Power-September-25-2026/README.md}{Manuscript page and citation}; \href{https://github.com/openai/math/blob/main/preprints/Foulkes-Conjecture-for-the-Sixth-Symmetric-Power-September-25-2026/main.pdf}{Latest version}; Statement source: \href{https://github.com/openai/math/blob/adc7f1241b42e322a6451854ab7e4b4c146bf78a/preprints/Foulkes-Conjecture-for-the-Sixth-Symmetric-Power-September-25-2026/build/sections/introduction.tex\#L19}{thm:main}.
\item Family 210: \href{https://github.com/openai/math/blob/adc7f1241b42e322a6451854ab7e4b4c146bf78a/lean/docs/210.md}{Lean scope}.
\item \href{https://github.com/openai/math/blob/adc7f1241b42e322a6451854ab7e4b4c146bf78a/CONTENTS.md}{OpenAI catalogue entry}.
\end{itemize}
\end{document}
